% !PRINTED 399
\weberpage{399}
Daraus folgt:

1. \gesperrt{Für ein negatives $D$ gibt es im allgemeinen nur die zwei Einheiten}
\[ \varepsilon = +1, \qquad \varepsilon = -1, \]
\gesperrt{und für zwei besondere Fälle}
\[
\begin{aligned}
D=-4\colon\quad &\varepsilon=+1,-1,+i,-i,\\
D=-3\colon\quad &\varepsilon=+1,-1,\\
&-\frac{1+\sqrt{-3}}{2},-\frac{1-\sqrt{-3}}{2},+\frac{1+\sqrt{-3}}{2},+\frac{1-\sqrt{-3}}{2}
\end{aligned}
\]

Die von Null verschiedenen Zahlen der Ordnung $O$ sind dann zu je zwei oder in den beiden Ausnahmefällen zu je vier oder zu je sechs assoziiert.

2. \gesperrt{Ist $D$ positiv, so gibt es eine Lösung der Pellschen Gleichung:}
\[ T^2 - U^2 D = \pm 4, \]
\gesperrt{in der $T$ und $U$ positive und möglichst kleine Werte haben. Ist dann}
\[ \varepsilon = \frac{T + U\sqrt{D}}{2}, \]
\gesperrt{so sind alle Einheiten in der Form $\pm \varepsilon^{\pm\nu}$ enthalten, wo $\nu$ die Reihe der Zahlen $0, 1, 2, \ldots$ durchläuft.}

Ist die Gleichung (2) für das negative Zeichen lösbar, so ist $N(\varepsilon) = -1$, und es hat $\pm \varepsilon^{\pm\nu}$ für gerade $\nu$ eine positive, für ungerade $\nu$ eine negative Norm. Ist aber (2) nur für das obere Zeichen lösbar, so haben alle Einheiten positive Norm. Unter den Zahlen $y$ der Ordnung $O$, die nicht Einheiten sind, gibt es aber immer solche mit negativer Norm.

Ist $N(\varepsilon) = -1$, so gibt es zu jeder Zahl $y$ in $O$ mit negativer Norm eine assoziierte mit positiver Norm. Ist aber $N(\varepsilon) = +1$, so zerfallen die Zahlen $y$ in zwei Klassen, von denen die eine positive und die andere negative Norm hat, und keine Zahl der einen Klasse ist mit einer der anderen Klasse assoziiert.

Demgemäß betrachten wir nur die Zahlen $y$ der Ordnung $O$ mit positiver Norm, unter denen $y$ und $\pm \varepsilon^{\pm\nu} y$ nur dann assoziiert sind, wenn
\[ \varepsilon = \frac{T + U\sqrt{D}}{2} \]
%CONT>>
